\SourceLeaf{408}%
\[
\text{(2)}\quad .\ .\ .\ .\ .\quad\cos b=\cos a.\cos c+\sin a.\sin c.\cos B
\]
Dividiert man (A) durch \(\sin c\), multipliziert (2) mit \(\operatorname{cotg}c\) und addiert, erhält man
\[
\frac{\cos a}{\sin c}=\frac{\cos a\cos c^2}{\sin c}+\sin b.\cos A+\sin a.\cos c.\cos B
\]
also
\[
\text{(3)}\quad .\ .\ .\quad\cos a.\sin c=\sin b.\cos A+\sin a.\cos c.\cos B
\]
und ebenso
\[
\text{(4)}\quad .\ .\ .\quad\cos c.\sin a=\sin b.\cos C+\sin c.\cos a.\cos B
\]
\par
Multipliziert man (4) mit \(\cos B\) und addiert (3), erhält man
\[
\text{(5)}\quad .\ .\quad\cos a.\sin c.\sin B^2=\sin b.\cos A+\sin b.\cos B.\cos C
\]
und ebenso
\[
\text{(6)}\quad .\ .\quad\cos a.\sin b.\sin C^2=\sin c.\cos A+\sin c.\cos B.\cos C
\]
\par
Multipliziert man (5) mit \(\frac{\sin c}{\cos a}\) und zieht davon (6), multipliziert mit \(\frac{\sin b}{\cos a}\), ab, so erhalten wir
\[
\text{(7)}\quad .\ .\ .\ .\ .\quad\sin c^2.\sin B^2-\sin b^2.\sin C^2=0
\]
oder
\[
\text{(B)}\quad .\ .\ .\ .\ .\ .\quad\sin c.\sin B=\sin b.\sin C
\]
Dividiert man (A,\,5) durch \(\sin b\) und zieht davon (B), multipliziert mit \(\frac{\cos a\sin B}{\sin b}\), ab, ergibt sich
\[
\text{(C)}\quad .\ .\ .\ .\quad\cos A=-\cos B.\cos C+\sin B.\sin C.\cos a
\]
Dividiert man schließlich (A,\,4) durch \(\sin c\) und zieht davon (B), multipliziert mit \(\frac{\operatorname{cotg}C}{\sin c}\), ab, ergibt sich
\[
\text{(D)}\quad .\ .\ .\ .\quad\operatorname{cotg}c.\sin a-\operatorname{cotg}C.\sin B=\cos a.\cos B
\]
\par
Aus diesen vier Grundformeln folgen unmittelbar die sogenannten \textsc{Neper}schen Analogien und die Vereinfachungen, die möglich sind, wenn das sphärische Dreieck rechtwinklig sein soll.

\{Man vergleiche diese Herleitung mit der von \textsc{Lagrange} im sechsten Heft des \textit{Journal de l’Ecole Polytechnique}.\}
